\chapter[Sermon 16. He Allows and Commends Those That Covet to Live Godly in the Church of Sardis, Exhorting Them That They Would So Hold on and Proceed.]{He Allows and Commends Those That Covet to Live Godly in the Church of Sardis, Exhorting Them That They Would So Hold on and Proceed.}
\label{sermon:016}
\markright{Sermon 16. He Allows and Commends Those That Covet to Live Godly ...}

But you have a few names in Sardis who have not defiled their garments; and they shall walk with me in white, for they are worthy. He who overcomes shall be clothed in white array, and I will not blot out his name from the book of life; and I will confess his name before my Father and before the angels. He who has ears let him hear what the Spirit says to the congregations.

The second part of this heavenly letter is contained in these points, wherein the innocence, holiness, and integrity of the faithful in the congregation of Sardis, in true religion, are praised and commended. He exhorts them with a very generous promise to persevere. Finally, he again proposes to them abundant rewards: even to the corrupt, if they repent; and to the faithful, if they continue as they are.

\index{Erasmus of Rotterdam}\index{Arethas of Caesarea}The Complutensian book reads thus: But you have a few names in Sardis. This is as much as to say that not all are corrupt and dead among you, although indeed those are very few. And so Arethas reads it in Greek, and the common translation in Latin, as found in other copies, which Erasmus follows: you have a few names also at Sardis; that is, even in Sardis you have names, but few. And he uses “names” to signify notable men. This manner of speaking is also in our language. For we say, “there is no man of name,” meaning no excellent or noble personage. He signifies therefore that there are in the same church noble personages, and that noble in soundness of faith and holiness of life, but very few, if they should be referred or compared to the number of hypocrites or the dead, which in truth are a great deal more. Neither ought we here to marvel. For the Lord says also in the Gospel that many are called, few are chosen; and that the greater part of this world walks in that broad and wide way of perdition, Matthew 20:7. Which also Saint Peter repeated in the second chapter of his second epistle. Those who seek to defend their error by a multitude are rather to be hissed at than confuted. You shall hear very often at this day: “We are but a few in number, we are innumerable, and therefore our matter is the better.”

But that same excellent thing is chiefly to be observed, that although they were but few good, yet nevertheless the Lord commends and extols those few, doubtless for the example and imitation of all other churches. The words in deed are short, but the praise most ample and large. That they had not defiled their garments: which is as much as if he had said, you have not polluted your souls with strange opinions or spots of heresy. For you have remained sincere in the true faith: your bodies also, and the whole conversation of your life, you have not defiled with filthy lusts, with fleshly pleasures and voluptuousness. Doubtless this is the greatest praise and most certain sign of perfect godliness: wherewith I would wish that more of us were marked. But the manner of speech here requires also an exposition. The allegory of garments is often and much used in holy scripture. The use of apparel invented of God himself, and showed to our forefathers, has this chief property, to hide the privy parts of our body, to beautify and set forth the body, and keep of heat and cold. And therefore Christ himself is called the garment of Christians, and in the gospel in deed the wedding garment. Whereupon the Apostle advises us to put on the new man, which is made after God even Christ himself. Romans 13, Ephesians 4, Colossians 3. For Christ covers not only our privy parts, but all the filthiness also of the soul, he adorns and beautifies us, and drives from us all injury, and all evil. And we defile this garment, when neither in faith nor in holiness of life we do answer to our profession. For Christ is our garment, and Christianity, sincere faith, and holiness of life are our apparel: and even faith and our conversation is our garment. Forasmuch therefore as the Sardinians were of a sincere faith, and incorrupt manners, they are said to have kept their garments clean and undefiled. The Lord also gives now a reward unto virtue. And they shall walk with me, says he, in white array. These excellent things verily doth he rehearse to retain the Sardinians in their duty, to nourish them to greater things and to move other also to sincerity and integrity. Saints walk with Christ in white array, that is to say, have fruition of the same glory, wherein we believe Christ to shine. For he desires his father, that he will grant to the faithful, that wherever he is, they may be with him, and see his glory, etc., in John 17. And with Saint Matthew, in the transformation or clarifying, the face of Christ appeared bright like the sun, his apparel and rest of his body as light. So appeared Christ unto John in the first chapter of this book, clothed in white array. Now therefore says he, the godly that have not defiled their garment, shall accompany me, having put on light also. He adds another thing, for they be worthy. This is the greatest praise, when the Captain says, that the soldier is worthy of honor and glory. The greatest shame or ignominy is, when it is said with us, thou art unworthy. The first kind of speech shows him to be most excellent in all kind of virtue, which is said to be worthy of eternal light; by the latter is signified, that he which is accounted unworthy of a good and excellent thing, is marvelous negligent and ungracious. But here we need not to reason of the merit and desert of worthiness. God pronounces his to be worthy of glory; the godly refer all the goodness that is in them unto grace, and still complain of their unworthiness. Not to reprove God of lying, but to praise and commend the excellent goodness that is in him, acknowledging in deed that he rewards good works, and dignifies the worthiness of saints: but they are nothing proud hereof, but acknowledge all this to come of grace. This appears in the doctrine of the Gospel, Luke 17, Matthew 25, where saints commended of God, for the works of mercy, seem to acknowledge nothing thereof.

\index{Jerome, Saint}\index{Repentance}However, he declares at length the most ample promises of God, whereby he may not only retain in their duty the saints and undefiled Sardinians, but might also reduce all others who go astray at all times into the way of repentance, integrity, and holiness. And three things he promises: first, indeed, white apparel, that is to say, glorifying and everlasting light, and the glorious company of Christ, of which I have spoken already. Secondly, and I will not put out his name out of the book of life. For just as cities have books wherein the names of their citizens are written, so too is God in the scriptures said to have, after the manner of men, a book of life, or of his elect. What that book is, and whose name is written therein, none of us can tell, since none has looked therein. We must learn from the scriptures who are the citizens of the kingdom of God. For that their names are written in the book of life, no one need doubt. And Jerome says: as many as have believed, he has given them power that they may be made the children of God. Paul says: He who does not have the spirit of Christ is none of his. And the spirit cries in the minds of the godly, “Abba, Father.” The same apostle says: God has predestined us that he might adopt us as his children through Jesus Christ. Moreover: he has chosen us in Christ before the foundations of the world were laid. Therefore, all believers are written in that celestial number. Whoever therefore does not believe, or does not persevere in the faith, either they are not written in the book of life, or else they are put out again from the book of life. Finally, the Son acknowledges the believers and those who persevere in the true faith before his heavenly Father and his angels. And here he repeats the evangelical doctrine from the 10th chapter of Matthew and the 8th of Mark. And it is undoubtedly a great matter in that universal judgment to be known by the Son of God, by the high judge, to be greeted and friendly spoken to by him, and that to our great praise. If any prince would in a great assembly of people know you, yes, embrace and commend you, how happy and fortunate would you think yourself? But then the very Son of God, king of kings, and lord of lords, will embrace you. Let us think of these things in time and amend our manners.

\index{John, Saint}For all these things pertain to us, the final and resounding declaration of Saint John confirms: “Let him who has ears, hear,” and so on, as we have discussed elsewhere. To the Lord be praise and glory.
